\documentclass[tikz]{standalone}

\usepackage{tikz}
\usepackage{tkz-euclide}
\usepackage{fourier-otf}
\usepackage{fontspec}

\usetikzlibrary{calc}
\usetikzlibrary{3d}

\tikzset{
  every picture/.append style={scale=2, every node/.style={scale=2}}
}



\begin{document}
  \begin{tikzpicture}[scale=0.375]
    % Configuration des limites
    \def\xmin{-8}
    \def\xmax{9}
    \def\ymin{-6}
    \def\ymax{8}
    
    % Axes avec flèches
    \draw[thick,->] (\xmin,0) -- (\xmax,0);
    \draw[thick,->] (0,\ymin) -- (0,\ymax);
    
    % Graduations sur les axes
    \foreach \x in {\xmin,...,-1,1,2,...,\xmax} {
        \draw (\x,0.1) -- (\x,-0.1);
    }
    \foreach \y in {\ymin,...,-1,1,2,...,\ymax} {
        \draw (0.1,\y) -- (-0.1,\y);
    }
    
    % Vecteurs unitaires
    \draw[thick,->] (0,0) -- (1,0);
    \node[below] at (0.5,0) {$\vec{\imath}$};
    \draw[thick,->] (0,0) -- (0,1);
    \node[left] at (0,0.5) {$\vec{\jmath}$};
    
    % Fonction f(x) = -1 + e^x (en rouge)
    \draw[red, thick, domain=\xmin:\xmax, samples=200, smooth] 
        plot (\x, {-1 + exp(\x)});
    
    % Fonction réciproque f^(-1)(x) = ln(x+1) (en vert)
    % Domaine: x > -1
    \draw[green, thick, domain=-0.99:\xmax, samples=200, smooth] 
        plot (\x, {ln(\x + 1)});
    
    % Bissectrice y = x (en bleu)
    \draw[blue, thick, domain=-5:8] 
        plot (\x, \x);
    
    % Étiquettes
    \node[below left] at (6,1.945) {$\mathcal{C}_{f}$};
    \node[below left] at (5,8) {$\mathcal{C}_{f^{-1}}$};
    \node[above right, rotate=45] at (6,6) {$y=x$};
        
\end{tikzpicture}
\end{document}
